\section*{Aktivitas Pembelajaran dan Praktikum}
\addcontentsline{toc}{section}{Aktivitas Pembelajaran dan Praktikum}

\begin{table}[H]
\centering
\caption{Rencana aktivitas Bab XIII}
\label{tab:ch13-aktivitas}
\begin{tabularx}{\textwidth}{@{}llX@{}}
\toprule
\textbf{Sesi} & \textbf{Bentuk} & \textbf{Kegiatan} \\
\midrule
1 & Ceramah & Arsitektur parser generator dan Bison \\
2 & Live coding & Menulis grammar dan semantic action \\
3 & Latihan & Membaca laporan konflik Bison \\
4 & Praktikum & M3: Bison + AST + integrasi \\
\bottomrule
\end{tabularx}
\end{table}

\subsection*{M3 / Praktikum 3: Bison dan AST}
\textbf{Tujuan.} Membangun parser Bison untuk Subset-C dan
menghasilkan AST lengkap.

\textbf{Keluaran.}
\begin{enumerate}[label=\arabic*.,leftmargin=*]
  \item Grammar Bison untuk ekspresi, pernyataan, dan deklarasi.
  \item Semantic action yang membangun AST.
  \item Integrasi Flex--Bison dan driver.
  \item Minimal 20 kasus uji: valid dan tidak valid.
  \item Perbandingan dengan recursive descent.
\end{enumerate}

\textbf{Langkah kerja.}
\begin{enumerate}[label=\arabic*.,leftmargin=*]
  \item Lengkapi \file{src/parser.y}.
  \item Definisikan \code{\%union}, \code{\%token}, dan \code{\%type}.
  \item Tulis semantic action untuk setiap aturan.
  \item Integrasikan \code{yylex} dari lexer.
  \item Jalankan Bison dan perbaiki konflik.
  \item Uji dan bandingkan hasilnya dengan parser tangan.
\end{enumerate}

\textbf{Rubrik.}
\begin{table}[H]
\centering
\caption{Rubrik Praktikum 3}
\label{tab:ch13-rubrik}
\begin{tabularx}{\textwidth}{@{}lXr@{}}
\toprule
\textbf{Aspek} & \textbf{Indikator} & \textbf{Bobot} \\
\midrule
Kebenaran & Grammar dan AST benar & 40\% \\
Kualitas kode & Aksi ringkas, penamaan, komentar & 30\% \\
Kelengkapan & Ekspresi, pernyataan, deklarasi & 20\% \\
Pengujian & Kasus normal, batas, kesalahan & 10\% \\
\bottomrule
\end{tabularx}
\end{table}

\textbf{Definisi selesai.}
\begin{checklist}
  \item Grammar Bison tanpa konflik tak terduga.
  \item AST terbentuk untuk program valid.
  \item 20 kasus uji tersedia.
  \item Perbandingan dengan recursive descent terdokumentasi.
\end{checklist}
